\chapter{Menú principal e Inicio}\label{cap:menu-inicio}

% ---------------------------------------------------------------------
\section{Menú principal}\label{sec:menu-principal}

\paraque Lleva a cada sección de Qadra. Cada rol ve solo las opciones que
puede usar.

\disponible Administrador, Litigante y Asistente legal. El menú del Cliente
se describe en la Parte B.

El menú está en el costado izquierdo de la pantalla. La \cref{tab:menu}
muestra sus opciones.

\begin{table}[H]
  \centering
  \caption{Opciones del menú principal por rol.}\label{tab:menu}
  \begin{tabularx}{\textwidth}{L{3.4cm} Y C{1.3cm} C{1.3cm} C{1.4cm}}
    \toprule
    \textbf{Opción} & \textbf{Qué abre} & \textbf{Adm.} & \textbf{Litig.} & \textbf{Asist.} \\
    \midrule
    \textbf{Inicio}            & Resumen del día y accesos rápidos. & Sí & Sí & Sí \\
    \textbf{Expedientes}       & Archivo de expedientes y cada expediente abierto. & Sí & Sí & Sí \\
    \textbf{Por firmar}        & Contratos que firma por el despacho. & Sí & Sí & No \\
    \textbf{Informes}          & Informes de expedientes en PDF y CSV. & Sí & Sí & No \\
    \textbf{Agenda}            & Audiencias y vencimientos de sus expedientes. & Sí & Sí & Sí \\
    \textbf{Alertas}           & Sus avisos personales y preferencias. & Sí & Sí & Sí \\
    \textbf{Equipo}            & Cuentas e invitaciones del despacho. & Sí & No & No \\
    \textbf{Bitácora}          & Movimientos e incidentes de integridad. & Sí & No & No \\
    \textbf{Calendarios}       & Calendarios jurisdiccionales. & Sí & No & No \\
    \textbf{Mi firma digital}  & Su clave, certificado y vínculo de firma. & Sí & Sí & No \\
    \textbf{Guía de uso}       & Ayuda breve dentro de Qadra. & Sí & Sí & Sí \\
    \bottomrule
  \end{tabularx}
\end{table}

\begin{itemize}
  \item \textbf{Equipo}, \textbf{Bitácora} y \textbf{Calendarios} aparecen
        bajo el título \textbf{Administración}.
  \item Para ganar espacio en pantallas grandes, haga clic en
        \textbf{Contraer menú}; para mostrarlo de nuevo, haga clic en
        \textbf{Expandir menú}.
  \item En el teléfono, el menú se abre con el botón \textbf{Abrir menú} de
        la barra superior.
\end{itemize}

\begin{figure}[H]
  \centering
  \capturapendiente{capturas/A-11-menu-principal.png}%
    {Pantalla completa con el menú lateral desplegado, incluida la sección Administración · Administrador}
  \caption{Menú principal del Administrador.}
  \label{fig:manual-menu-principal}
\end{figure}

% ---------------------------------------------------------------------
\section{Barra superior}\label{sec:barra-superior}

\paraque Ofrece la búsqueda rápida de expedientes, sus avisos sin leer y el
menú de su cuenta.

\disponible Administrador, Litigante y Asistente legal.

\begin{table}[H]
  \centering
  \caption{Elementos de la barra superior.}\label{tab:barra-superior}
  \begin{tabularx}{\textwidth}{L{3.6cm} Y}
    \toprule
    \textbf{Elemento} & \textbf{Cómo se usa} \\
    \midrule
    \textbf{Buscar expedientes}
      & Escriba al menos 2 caracteres del título, la referencia, el NUC o la
        carpeta judicial. Seleccione un resultado para abrirlo. También se
        activa con \texttt{Ctrl}+\texttt{K}. \\
    Campana de avisos
      & Muestra el panel \textbf{Avisos sin leer}. Haga clic en
        \textbf{Abrir la bandeja de alertas} para ver todos en
        \textbf{Alertas}. \\
    Círculo con sus iniciales (\textbf{Tu cuenta})
      & Muestra su nombre, correo y rol, y el botón
        \textbf{Cerrar sesión} (\cref{sec:cerrar-sesion}). \\
    \bottomrule
  \end{tabularx}
\end{table}

\nota{La búsqueda rápida cubre los primeros 300 expedientes. Si no encuentra
uno, haga clic en \textbf{Ir al archivo de expedientes} y use los filtros
(\cref{sec:buscar-expediente}).}

% ---------------------------------------------------------------------
\section{Inicio}\label{sec:inicio}

\paraque Muestra su situación del día y lo lleva rápido a sus pendientes.

\disponible Administrador, Litigante y Asistente legal. El contenido cambia
según el rol.

\subsection*{Inicio del Administrador y del Litigante}

El Administrador ve los datos de \textbf{Todo el despacho}; el Litigante, los
de \textbf{Tus expedientes asignados}.

\begin{itemize}
  \item \textbf{Indicadores:} \textbf{Expedientes activos},
        \textbf{Plazos vencidos}, \textbf{Vencen en 7 días},
        \textbf{Por revisar} y \textbf{Contratos en firma}. Haga clic en
        \textbf{Actualizar} para recalcularlos.
  \item \textbf{Próximos plazos y audiencias:} actividades de los próximos
        30 días. Haga clic en una actividad o en \textbf{Ver agenda} para
        abrir la \textbf{Agenda}.
  \item \textbf{Carga por litigante:} cuántos expedientes activos tiene
        asignados cada litigante.
\end{itemize}

\begin{figure}[H]
  \centering
  \capturapendiente{capturas/A-12-inicio-tablero.png}%
    {Inicio con indicadores, próximos plazos y audiencias y carga por litigante · Litigante}
  \caption{Inicio del Administrador y del Litigante.}
  \label{fig:manual-inicio-tablero}
\end{figure}

\subsection*{Inicio del Asistente legal}

\begin{itemize}
  \item \textbf{Tu expediente actual:} haga clic en
        \textbf{Abrir expediente} para continuar con el último expediente,
        o en \textbf{Seleccionar expediente} si aún no eligió uno.
  \item \textbf{Accesos:} atajos a \textbf{Expedientes},
        \textbf{Pendientes de sello}, \textbf{Documentos sellados} y
        \textbf{Mis alertas}.
  \item \textbf{Tu siguiente paso:} recuerda el recorrido de un documento:
        \textbf{Carga}, \textbf{Sella} y \textbf{Verifica y descarga}.
\end{itemize}

\begin{figure}[H]
  \centering
  \capturapendiente{capturas/A-13-inicio-asistente.png}%
    {Inicio con Tu expediente actual, Accesos y Tu siguiente paso · Asistente legal}
  \caption{Inicio del Asistente legal.}
  \label{fig:manual-inicio-asistente}
\end{figure}

\resultado Desde \textbf{Inicio} llega en un clic a sus expedientes, su
agenda o sus documentos pendientes.

\nota{La opción \textbf{Guía de uso} del menú resume, dentro de Qadra, las
tareas más comunes y responde preguntas frecuentes.}
